% GENERATED FILE. Edit canonical lessons, then run npm run generate:books.
% canonical-source-hash: fnv1a64:4520d2021a4073e8
% canonical-lessons: TE-C140-suraksitam, TE-C140-abaddham, TE-C140-budida, TE-C140-bangaru, TE-C140-padaina

\chapter{Safe, Lie, Grey, Golden, Spoiled}
\label{ch:te-safe-lie-grey-golden-spoiled}
\hlchaptermodality{140}

\begin{chapteropening}
\textbf{By the end of this chapter:} \emph{I can say safe, a lie, grey, golden and spoiled.}

Complete the last lesson of chapter 140: I can say safe, a lie, grey, golden and spoiled.
\end{chapteropening}

\section[\texorpdfstring{\te{సురక్షితం} (surakṣitaṁ)}{సురక్షితం (surakṣitaṁ)}]{\te{సురక్షితం} (surakṣitaṁ) — safe}

\label{lesson:TE-C140-suraksitam}



\begin{quote}
Before the new one: say the Telugu for old (of people), then the Telugu for famous.
\end{quote}

\subsection*{You'll want to know: \te{సురక్షితం}}
\textbf{\te{సురక్షితం}} — \emph{surakṣitaṁ} — \textquotedblleft{}safe\textquotedblright{}.

\begin{cousinweb}
\textbf{\te{సురక్షితం}} (\emph{surakṣitaṁ}) is Sanskrit \emph{surakṣita}, \textquotedblleft{}well protected\textquotedblright{}, taken over whole as a learned loan, a tatsama, with the Telugu noun ending \emph{-ṁ} (CDIAL 13499).
\end{cousinweb}

\begin{grammarlens}[title={What you've built}]
One more word for this chapter.
\end{grammarlens}

\subsection*{Your turn}
\begin{itemize}
\raggedright
  \item \emph{Say it:} \emph{surakṣitaṁ}
  \item \emph{Say it:} \emph{surakṣitaṁ}, once more
  \item \emph{Say it:} say \emph{surakṣitaṁ}
  \item \emph{Recall:} say the Telugu for old (of people), then the Telugu for famous, then say \emph{surakṣitaṁ} again
\end{itemize}

\begin{tcolorbox}[breakable,colback=teal!4,colframe=teal!35!black,title={Before you move on}]
What does \te{సురక్షితం} mean? (\textquotedblleft{}Safe\textquotedblright{}.) Say it once more.
\end{tcolorbox}
\section[\texorpdfstring{\te{అబద్ధం} (abaddhaṁ)}{అబద్ధం (abaddhaṁ)}]{\te{అబద్ధం} (abaddhaṁ) — a lie; untrue}

\label{lesson:TE-C140-abaddham}



\begin{quote}
Before the new one: say the Telugu for famous, then the Telugu for safe.
\end{quote}

\subsection*{You'll want to know: \te{అబద్ధం}}
\textbf{\te{అబద్ధం}} — \emph{abaddhaṁ} — \textquotedblleft{}a lie; untrue\textquotedblright{}.

\begin{cousinweb}
\textbf{\te{అబద్ధం}} (\emph{abaddhaṁ}) is Sanskrit \emph{abaddha}, \textquotedblleft{}unconnected, meaningless\textquotedblright{}, taken over whole as a learned loan, a tatsama. Telugu uses it for a lie.
\end{cousinweb}

\begin{grammarlens}[title={What you've built}]
One more word for this chapter.
\end{grammarlens}

\subsection*{Your turn}
\begin{itemize}
\raggedright
  \item \emph{Say it:} \emph{abaddhaṁ}
  \item \emph{Say it:} \emph{abaddhaṁ}, once more
  \item \emph{Say it:} say \emph{abaddhaṁ}
  \item \emph{Recall:} say the Telugu for famous, then the Telugu for safe, then say \emph{abaddhaṁ} again
\end{itemize}

\begin{tcolorbox}[breakable,colback=teal!4,colframe=teal!35!black,title={Before you move on}]
What does \te{అబద్ధం} mean? (\textquotedblleft{}A lie; untrue\textquotedblright{}.) Say it once more.
\end{tcolorbox}
\section[\texorpdfstring{\te{బూడిద} (būḍida)}{బూడిద (būḍida)}]{\te{బూడిద} (būḍida) — grey; ash}

\label{lesson:TE-C140-budida}



\begin{quote}
Before the new one: say the Telugu for safe, then the Telugu for a lie; untrue.
\end{quote}

\subsection*{You'll want to know: \te{బూడిద}}
\textbf{\te{బూడిద}} — \emph{būḍida} — \textquotedblleft{}grey; ash\textquotedblright{}.

\begin{cousinweb}
\textbf{\te{బూడిద}} (\emph{būḍida}) is native Dravidian. DEDR 4316 lists \emph{būḍida}, \textquotedblleft{}ashes\textquotedblright{}, with Tamil \emph{puḻuti}, \textquotedblleft{}dust\textquotedblright{}. DEDR compares Sanskrit \emph{bhūti}, \textquotedblleft{}ashes\textquotedblright{}, and notes that some forms in the entry are probably borrowed from it.
\end{cousinweb}

\begin{grammarlens}[title={What you've built}]
One more word for this chapter.
\end{grammarlens}

\subsection*{Your turn}
\begin{itemize}
\raggedright
  \item \emph{Say it:} \emph{būḍida}
  \item \emph{Say it:} \emph{būḍida}, once more
  \item \emph{Say it:} say \emph{būḍida}
  \item \emph{Recall:} say the Telugu for safe, then the Telugu for a lie; untrue, then say \emph{būḍida} again
\end{itemize}

\begin{tcolorbox}[breakable,colback=teal!4,colframe=teal!35!black,title={Before you move on}]
What does \te{బూడిద} mean? (\textquotedblleft{}Grey; ash\textquotedblright{}.) Say it once more.
\end{tcolorbox}
\section[\texorpdfstring{\te{బంగారు} (baṅgāru)}{బంగారు (baṅgāru)}]{\te{బంగారు} (baṅgāru) — golden}

\label{lesson:TE-C140-bangaru}



\begin{quote}
Before the new one: say the Telugu for a lie; untrue, then the Telugu for grey; ash.
\end{quote}

\subsection*{You'll want to know: \te{బంగారు}}
\textbf{\te{బంగారు}} — \emph{baṅgāru} — \textquotedblleft{}golden\textquotedblright{}.

\begin{cousinweb}
The origin of \textbf{\te{బంగారు}} (\emph{baṅgāru}) is uncertain. Neither DEDR nor CDIAL records it.
\end{cousinweb}

\begin{grammarlens}[title={What you've built}]
One more word for this chapter.
\end{grammarlens}

\subsection*{Your turn}
\begin{itemize}
\raggedright
  \item \emph{Say it:} \emph{baṅgāru}
  \item \emph{Say it:} \emph{baṅgāru}, once more
  \item \emph{Say it:} say \emph{baṅgāru}
  \item \emph{Recall:} say the Telugu for a lie; untrue, then the Telugu for grey; ash, then say \emph{baṅgāru} again
\end{itemize}

\begin{tcolorbox}[breakable,colback=teal!4,colframe=teal!35!black,title={Before you move on}]
What does \te{బంగారు} mean? (\textquotedblleft{}Golden\textquotedblright{}.) Say it once more.
\end{tcolorbox}
\section[\texorpdfstring{\te{పాడైన} (pāḍaina)}{పాడైన (pāḍaina)}]{\te{పాడైన} (pāḍaina) — spoiled}

\label{lesson:TE-C140-padaina}



\begin{quote}
Before the new one: say the Telugu for grey; ash, then the Telugu for golden.
\end{quote}

\subsection*{You'll want to know: \te{పాడైన}}
\textbf{\te{పాడైన}} — \emph{pāḍaina} — \textquotedblleft{}spoiled\textquotedblright{}.

\begin{cousinweb}
\textbf{\te{పాడైన}} (\emph{pāḍaina}) is native Dravidian: \emph{pāḍu}, \textquotedblleft{}ruin\textquotedblright{}, one word with Tamil \emph{pāḻ} (DEDR 4110), and \emph{ayina}, \textquotedblleft{}which has become\textquotedblright{} (DEDR 333). DEDR compares Prakrit \emph{pāla}, \textquotedblleft{}worn out\textquotedblright{}.
\end{cousinweb}

\begin{grammarlens}[title={What you've built}]
One more word for this chapter.
\end{grammarlens}

\subsection*{Your turn}
\begin{itemize}
\raggedright
  \item \emph{Say it:} \emph{pāḍaina}
  \item \emph{Say it:} \emph{pāḍaina}, once more
  \item \emph{Say it:} say \emph{pāḍaina}
  \item \emph{Recall:} say the Telugu for grey; ash, then the Telugu for golden, then say \emph{pāḍaina} again
\end{itemize}

\begin{tcolorbox}[breakable,colback=teal!4,colframe=teal!35!black,title={Before you move on}]
What does \te{పాడైన} mean? (\textquotedblleft{}Spoiled\textquotedblright{}.) Say it once more.
\end{tcolorbox}
